\documentclass[11pt]{article}

% ===============================
% LuaLaTeX setup
% ===============================
\usepackage{fontspec}
\setmainfont{TeX Gyre Pagella}
\setsansfont{TeX Gyre Heros}
\setmonofont{TeX Gyre Cursor}

% ===============================
% Mathematics and layout
% ===============================
\usepackage{amsmath, amssymb, mathtools}
\usepackage{geometry}
\geometry{margin=1in}
\usepackage{setspace}
\onehalfspacing

\usepackage{hyperref}
\usepackage{microtype}

% ===============================
% Metadata
% ===============================
\title{Irreversible Dynamics of Entropic Manifolds}
\author{Flyxion}
\date{\today}

\begin{document}
\maketitle

\begin{abstract}
This essay develops a unified architectural perspective on a family of entropic and geometric frameworks for modeling physical, informational, and semantic systems under constraint. Rather than proposing an additional fundamental ontology, the work synthesizes scalar entropy fields, vectorial flow structures, and geometric constraint formalisms into a single descriptive language. The central methodological commitment is constraint-first reasoning: admissible structure is specified prior to dynamics, and evolution arises as a secondary consequence of compatibility, conservation, and dissipation bounds. This perspective replaces expansionist and substance-centered ontologies with a view grounded in smoothing, descent, and relational flow, offering a common explanatory substrate across physics, cosmology, and cognition.
\end{abstract}

\section{Motivation: From Objects to Fields of Constraint}

Across contemporary theoretical physics, information theory, and cognitive modeling, explanatory frameworks remain implicitly object-centered. Systems are typically described as collections of entities endowed with intrinsic properties, acted upon by forces, and embedded in backgrounds that are taken as ontologically primitive. Constraint, dissipation, and context are then introduced as secondary corrections to otherwise free dynamics. While this methodology has proven empirically powerful, it exhibits persistent conceptual strain when confronted with non-equilibrium phenomena, emergent structure, and cross-domain unification (Jaynes, 1957; Prigogine, 1980).

The framework developed here reverses this hierarchy. Constraint is treated as primitive, while objects, forces, and even geometry are understood as stabilized patterns within constrained fields of variation. What exists, on this view, is not primarily substance but admissible transformation. Fields do not represent material stuff distributed over space; rather, they encode the degrees of freedom locally available without violating global or local constraints. Structure is what remains when entropy is allowed to flow, but not arbitrarily.

This inversion aligns naturally with constraint-based interpretations of statistical mechanics and thermodynamics, in which macroscopic regularities arise from admissibility conditions rather than microscopic laws (Jaynes, 2003; Callen, 1985). It also resonates with philosophical accounts of emergence that treat higher-level structure as stabilized under compatibility relations rather than reducible to constituent dynamics (Butterfield, 2014).

Importantly, this approach does not deny the utility of object-based descriptions. Instead, it treats them as derived coordinate systems on a deeper space of constraint-respecting configurations. Objects persist because entropy flow can support them; they dissolve when it cannot. Geometry, thermodynamics, and information thus appear not as independent explanatory domains, but as complementary projections of a single ordering principle.

\subsection*{Formal Perspective}

At a schematic level, the constraint-first approach replaces the question
\emph{``What equations govern motion?''}
with
\emph{``Which transformations are admissible?''}
Dynamics then emerge as consistency conditions on admissible histories rather than as axiomatic prescriptions.

\section{Scalar Entropy Fields and Vectorial Flow}

At the core of the integrated framework lies a scalar field $\Phi$, defined over a base manifold $M$, and interpreted as a measure of locally available degrees of freedom compatible with constraint. In physical contexts, $\Phi$ is naturally identified with entropy density; in informational and cognitive contexts, it corresponds to uncertainty or underdetermination; and in geometric contexts, it reflects smoothing capacity. These interpretations are not competing meanings but domain-specific realizations of a single abstract role (Cover and Thomas, 2006).

The base manifold $M$ itself is not assumed to be an ontological primitive. It represents the weakest differentiable structure required to index admissible variation. In physics, empirical constraints select spacetime as this structure; in cognition, the manifold corresponds to a representational state space. More generally, $M$ should be understood as an atlas induced by constraint coherence rather than as a pre-existing container.

Coupled to the scalar field $\Phi$ is a vector field $\vec{v}$ representing directed transport of admissible degrees of freedom. In thermodynamic systems, $\vec{v}$ may encode energy or mass flux; in informational systems, it may represent inferential propagation or attention flow. The defining role of $\vec{v}$ is to mediate redistribution of $\Phi$ across $M$.

The interaction between $\Phi$ and $\vec{v}$ is governed by a continuity-type relation
\[
\partial_t \Phi + \nabla \cdot (\Phi \vec{v}) = \sigma,
\]
where $\sigma$ denotes entropy production density. Crucially, admissibility requires $\sigma \ge 0$ almost everywhere, expressing a generalized second-law condition without appealing to a primitive temporal arrow (Prigogine, 1980; Callen, 1985).

Flow arises in response to gradients in $\Phi$, while transport reshapes those gradients in turn. Stable structures emerge precisely where entropy descent, flux divergence, and production reach dynamic balance. No assumption of spatial expansion, external forcing, or privileged coordinates is required. Redistribution, not growth, is the fundamental mechanism.

\subsection*{Formal Derivation: Entropy Balance}

Let $\Omega \subset M$ be a compact region with boundary $\partial \Omega$. Integrating the continuity relation yields
\[
\frac{d}{dt} \int_\Omega \Phi \, dV
=
- \int_{\partial \Omega} \Phi \vec{v} \cdot \hat{n} \, dS
+ \int_\Omega \sigma \, dV.
\]
In closed regions, where boundary flux vanishes, total entropy is non-decreasing. In open systems, local decreases in $\Phi$ are permitted provided they are compensated by boundary flux or external dissipation. Structure formation thus appears as an accounting phenomenon: local stabilization is paid for elsewhere in the system.

The condition $\sigma \ge 0$ functions not as a dynamical equation but as a geometric admissibility constraint on histories. Just as causal structure restricts admissible curves in spacetime, entropy production positivity restricts admissible transformations in configuration space.

\bigskip

\section{Variational and Geometric Encoding of Dynamics}

Although the framework is explicitly constraint-first, it admits a precise variational and geometric formulation once admissible configurations have been specified. The purpose of variational structure here is not to elevate an action functional to ontological primacy, but to provide a compact representation of compatibility conditions governing entropy redistribution. Dynamics arise as consequences of admissibility, not as axioms imposed in advance.

Fields such as $\Phi$ and $\vec{v}$ are treated as sections of structured bundles over the base manifold $M$, with admissible variations encoded in their tangent and cotangent data. Geometry, in this setting, serves as a bookkeeping device: it records which transformations preserve constraints and which violate them. Singularities and degeneracies are not pathologies but structured indicators of transitions between constraint regimes (Gromov, 1999).

A natural variational representation is obtained by introducing an entropy-weighted action functional of the schematic form
\[
\mathcal{S}[\Phi,\vec{v}]
=
\int dt \int_M
\left(
\frac{1}{2} \Phi \lVert \vec{v} \rVert^2
-
U(\Phi)
\right) dV
-
\int dt \int_M
\lambda \left(
\partial_t \Phi + \nabla \cdot (\Phi \vec{v}) - \sigma
\right) dV,
\]
where $U(\Phi)$ is a convex potential encoding constraint stiffness, $\lambda$ is a Lagrange multiplier enforcing entropy balance, and $\sigma \ge 0$ represents admissible entropy production. The functional form of $U$ is not fixed universally; it is part of the constraint architecture that distinguishes domains and regimes.

Variation with respect to $\vec{v}$ yields
\[
\Phi \vec{v} = \Phi \nabla \lambda,
\]
implying that admissible transport is necessarily gradient-driven. Entropy flows along directions of steepest constraint relaxation, rather than along arbitrary vector fields. Variation with respect to $\Phi$ gives
\[
\frac{1}{2} \lVert \vec{v} \rVert^2
-
U'(\Phi)
-
\partial_t \lambda
-
\vec{v} \cdot \nabla \lambda
=
0,
\]
which may be interpreted as a generalized Hamilton--Jacobi relation for the entropy potential.

These relations do not define independent equations of motion. Rather, they encode mutual consistency between entropy balance, transport admissibility, and constraint pressure. Apparent dynamics emerge as the projection of these compatibility conditions onto the chosen coordinates.

\subsection*{Hamiltonian Formulation}

A Hamiltonian description follows by identifying the conjugate momentum
\[
\pi := \lambda,
\]
so that $(\Phi,\pi)$ form a canonical pair. The corresponding Hamiltonian functional is
\[
\mathcal{H}[\Phi,\pi]
=
\int_M
\left(
\frac{1}{2} \Phi \lVert \nabla \pi \rVert^2 + U(\Phi)
\right) dV,
\]
supplemented by the entropy balance constraint. Hamilton’s equations reproduce the continuity equation for $\Phi$ and the gradient-flow condition for $\vec{v}$.

Importantly, dissipation enters not through explicit friction terms but through the admissibility condition $\sigma \ge 0$, which restricts the class of allowable trajectories in phase space. The symplectic structure underlying the Hamiltonian formulation is therefore not globally preserved. Instead, entropy production induces a controlled contraction of admissible phase-space volume, rendering conservative Hamiltonian flow a singular limiting case (Arnold, 1989; Baez and Biamonte, 2018).

From this perspective, dissipation is not an external modification of conservative dynamics but a geometric feature of the constraint surface itself. The Hamiltonian formalism remains valid as an organizing scaffold even in non-equilibrium regimes, provided it is interpreted as encoding admissibility rather than reversibility.

\section{Constraint-First Methodology}

A defining methodological feature of the integrated framework is its refusal to begin with equations of motion. Instead, analysis proceeds by specifying admissible configurations, equivalence relations under which descriptions are redundant, conservation or boundedness conditions, and allowable rates of entropy production. Only once these elements are fixed do evolution equations appear, as necessary conditions for internal consistency.

This inversion of the traditional methodological order explains why apparently disparate models often converge to similar macroscopic behavior despite differing microscopic assumptions. Such convergence reflects shared constraint architectures rather than shared dynamics. What appears as universality is, on this view, the expression of deep compatibility conditions imposed by entropy flow under constraint (Jaynes, 2003).

Constraint modification is permitted, but not freely. Deformation of the constraint landscape—for example through changes in the potential $U(\Phi)$—is itself an irreversible process subject to entropy production bounds. Constraints may relax or harden only along admissible entropy-respecting trajectories. This second-order restriction prevents unconstrained flexibility while avoiding infinite regress: constraints evolve, but only at a cost (Sklar, 1993).

\subsection*{Formal Remark: Admissible Histories}

Let $\mathcal{H}$ denote the space of field histories $(\Phi(t),\vec{v}(t))$. The entropy production inequality $\sigma \ge 0$ induces a partial order on $\mathcal{H}$ by restricting admissible paths. Irreversibility is thus encoded as an order-theoretic property of histories rather than as a primitive temporal asymmetry. Time serves as an ordering parameter, not as an explanatory primitive.

\bigskip

\section{Entropy Production and Irreversibility}

A central advantage of the constraint-first perspective is that irreversibility need not be postulated as a primitive temporal asymmetry. Instead, irreversibility emerges as a structural property of admissible entropy flow. The arrow of time is not imposed as an external parameter but is derived from the geometry of entropy production constraints (Prigogine, 1980; Callen, 1985).

Let $\Phi(x,t)$ denote the scalar entropy field on the base manifold $M$, and let $\vec{v}(x,t)$ be the associated transport field. The local balance of entropy is governed by
\[
\partial_t \Phi + \nabla \cdot (\Phi \vec{v}) = \sigma,
\]
where $\sigma(x,t)$ denotes entropy production density. The admissibility condition
\[
\sigma \ge 0
\]
is not interpreted as a dynamical law but as a geometric restriction on allowable histories. Any transformation requiring $\sigma < 0$ would increase locally available degrees of freedom without compensating constraint relaxation and is therefore excluded by construction.

Integrating this relation over a compact region $\Omega \subset M$ yields
\[
\frac{d}{dt} \int_\Omega \Phi \, dV
=
- \int_{\partial \Omega} \Phi \vec{v} \cdot \hat{n} \, dS
+ \int_\Omega \sigma \, dV.
\]
In closed regions, total entropy is non-decreasing. In open systems, local decreases in entropy are permitted only when balanced by boundary flux or external dissipation. Irreversibility thus appears as an accounting constraint rather than a metaphysical arrow.

Reversible evolution corresponds to the singular case $\sigma \equiv 0$. Such trajectories exist only on measure-zero subsets of admissible configuration space and are structurally unstable under perturbation. This instability explains the empirical rarity of reversible macroscopic processes without appealing to probabilistic typicality or microscopic reversibility arguments.

Importantly, entropy production does not imply disorder. On the contrary, localized increases in $\sigma$ can accompany the emergence of stable, low-entropy structures elsewhere, provided the global constraint architecture is respected. Irreversibility and structure formation are therefore mutually enabling rather than opposed.

\subsection*{Formal Remark: Entropy as Structural Asymmetry}

The condition $\sigma \ge 0$ induces a partial order on admissible histories. Temporal directionality arises from this ordering, not from the parameter $t$ itself. Time serves as a coordinate labeling admissible transformations; irreversibility is encoded in the shape of the admissible region of configuration space.

\section{Stability, Attractors, and Structure Formation}

Within the entropy–flow framework, the emergence of structure is a direct consequence of stability properties intrinsic to constrained entropy descent. Objects, patterns, and persistent configurations arise as attractors in the space of admissible fields, characterized by balanced entropy production and transport (Prigogine and Stengers, 1984; Butterfield, 2014).

Let $(\Phi_0,\vec{v}_0)$ denote a stationary configuration satisfying
\[
\partial_t \Phi_0 = 0,
\qquad
\nabla \cdot (\Phi_0 \vec{v}_0) = \sigma_0,
\qquad
\sigma_0 \ge 0.
\]
Consider perturbations of the form
\[
\Phi = \Phi_0 + \varepsilon \, \delta \Phi,
\qquad
\vec{v} = \vec{v}_0 + \varepsilon \, \delta \vec{v},
\]
and linearize the entropy balance equation to first order in $\varepsilon$. The resulting system defines a constrained linear operator whose spectrum determines the response of the configuration to admissible perturbations.

A configuration is said to be structurally stable if all admissible perturbations decay under entropy descent, modulo gauge equivalences. Stability is evaluated not with respect to arbitrary variations, but only those consistent with entropy production bounds and conservation constraints. This restriction substantially alters the classification of stable states relative to unconstrained dynamical systems.

Attractors in this framework are defined as compact subsets of field configuration space toward which entropy flow converges under admissible evolution. These attractors need not be fixed points; they may take the form of limit cycles, quasi-stationary manifolds, or slowly drifting configurations, provided entropy production remains locally balanced. What distinguishes an attractor is not immobility, but persistence under constraint.

From this perspective, objects are metastable regions of reduced entropy gradient, maintained by compensating fluxes in their surroundings. Apparent boundaries correspond to sharp transitions in constraint stiffness rather than to primitive surfaces. The interior of an object is characterized by suppressed entropy production relative to its environment, while the exterior supplies the dissipative cost required for its maintenance.

\subsection*{Formal Derivation: Linearized Stability}

Let $\mathcal{L}$ denote the linearized operator governing perturbations $(\delta \Phi, \delta \vec{v})$. Structural stability requires the spectrum of $\mathcal{L}$, restricted to admissible perturbations, to lie in the left half-plane. Eigenmodes corresponding to gauge directions are factored out. Nonlinear terms determine basin size and robustness, explaining why some admissible structures fail to appear under generic initial conditions.

Structure formation thus emerges as a generic outcome of constrained entropy descent. Complexity does not resist entropy; it is sculpted by it. Persistent structure is the stable residue left by irreversible flow under nontrivial constraints.

\bigskip

\section{Gauge, Symmetry, and Redundancy}

A constraint-first framework requires a principled distinction between physically or semantically meaningful variation and mere representational freedom. In the entropy–flow setting, this distinction is formalized through gauge structure. Gauge symmetries do not represent dynamical freedoms but encode equivalence relations between descriptions that preserve all admissible constraints (Baez and Pollard, 2017).

Let $\mathcal{F}$ denote the space of admissible field configurations $(\Phi,\vec{v})$ satisfying entropy balance and production bounds. A gauge transformation is defined as an automorphism of $\mathcal{F}$ that leaves invariant all observable entropy fluxes and production rates. Under such transformations, the numerical form of $\Phi$ and $\vec{v}$ may change, but all quantities of the form
\[
\int_\Omega \Phi \, dV,
\qquad
\int_{\partial \Omega} \Phi \vec{v} \cdot \hat{n} \, dS,
\qquad
\int_\Omega \sigma \, dV,
\]
remain unchanged for all admissible regions $\Omega \subset M$. Physical and semantic content is therefore identified not with individual field configurations, but with equivalence classes under the gauge group.

This reframing places symmetry on an operational footing. Symmetries are not postulated as fundamental invariances of nature, but arise as redundancies induced by constraint satisfaction. A continuous symmetry exists precisely when distinguishable change along a direction in configuration space is forbidden by admissibility. Conservation laws then follow as statements of identifiability rather than as metaphysical assertions of persistence (Gromov, 1999).

The presence of gauge structure clarifies the status of apparent degrees of freedom. Many variables that appear dynamical in unconstrained formulations are revealed, upon enforcing admissibility, to be pure gauge. Their apparent evolution corresponds to motion along equivalence classes rather than to genuine redistribution of entropy. Eliminating such redundancies simplifies the theory without loss of explanatory power.

In semantic and cognitive applications, gauge equivalence acquires a particularly transparent interpretation. Distinct representational encodings, narratives, or inferential pathways may be gauge-equivalent if they induce identical entropy flows and uncertainty reductions. Gauge equivalence does not imply equivalence of computational cost, accessibility, or robustness; it concerns invariance of constraint-respecting structure, not implementation efficiency. This distinction allows semantic flexibility without collapsing meaningful differences.

\subsection*{Formal Remark: Gauge Reduction}

Let $\mathcal{F}/\mathcal{G}$ denote the quotient of admissible configurations by the gauge group $\mathcal{G}$. Physical and semantic observables are functions on $\mathcal{F}/\mathcal{G}$. Evolution within a gauge orbit is unobservable; only motion transverse to gauge directions corresponds to genuine entropy redistribution.

\section{Cosmology Without Expansion}

When applied at cosmological scales, the entropy–flow framework offers an alternative to expansion-based interpretations of large-scale structure. Rather than treating metric expansion as a primitive mechanism, cosmological evolution is described as long-range redistribution of entropy within a constrained plenum. Space does not grow; it relaxes.

In this view, large-scale structure forms not because distances between points increase, but because entropy gradients smooth under constraint. Regions of high constraint pressure relax, driving redistribution that observers embedded within the flow interpret as kinematic expansion. Apparent horizons emerge as thermodynamic boundaries: surfaces beyond which entropy production along connecting paths renders information transfer effectively incoherent (Ellis, 2016).

Phenomena typically attributed to dark energy or accelerated expansion arise naturally as second-order effects of gradient relaxation. As entropy redistributes, local calibrations of clocks and rulers drift relative to distant structures, producing the observational signature of acceleration without invoking additional substances or modified gravitational dynamics.

This reinterpretation preserves the empirical successes of standard cosmology while altering its ontological commitments. Expansion becomes a derived description of entropic smoothing rather than a fundamental feature of spacetime. The universe is not an inflating container but a constrained distribution evolving toward smoother configurations.

\subsection*{Formal Perspective: Entropic Relaxation}

Let $\Phi$ be defined on cosmological scales. Long-wavelength modes dominate the entropy landscape, and their relaxation timescales set effective cosmological parameters. Apparent acceleration corresponds to the curvature of entropy descent trajectories in observational coordinates. Horizons correspond to regions where entropy production along null-like paths diverges, suppressing reconstructible signal transmission.

\bigskip

\section{Cosmological Observables and Entropic Smoothing}

When interpreted through the entropy–flow framework, standard cosmological observables admit a reformulation that does not rely on metric expansion as a primitive explanatory device. Instead, observed kinematic effects arise from cumulative entropy redistribution along observational histories. Measurements of distance, time, and energy are themselves constrained processes and therefore encode the global structure of entropy flow.

Let $M$ be a cosmological-scale manifold equipped with entropy field $\Phi$ and transport field $\vec{v}$. Consider a signal propagating along a curve $\gamma$ from emission event $e$ to observation event $o$. The observed redshift may be expressed schematically as
\[
1 + z
=
\exp\!\left(
\int_\gamma \alpha(\Phi, \nabla \Phi, \vec{v}) \, ds
\right),
\]
where $\alpha$ is a local attenuation functional determined by entropy gradients and flow divergence. This expression preserves the empirical scaling relations of standard cosmology while relocating their causal origin from spacetime expansion to cumulative entropic smoothing (Villani, 2009; Otto, 2001).

Cosmological horizons arise as thermodynamic boundaries rather than geometric singularities. Regions separated by sufficiently steep entropy gradients become effectively decoupled, not because signals cannot traverse expanding space, but because entropy production along connecting paths renders information transfer asymptotically incoherent. Horizons thus correspond to surfaces of maximal admissible dissipation.

This perspective yields principled empirical constraints. Any observation requiring sustained negative entropy production along causal histories, or demanding global entropy decrease without compensating boundary flux, would violate the framework. Likewise, phenomena irreducible to cumulative entropy-descent effects would resist reinterpretation within this architecture. The framework is therefore flexible but not unfalsifiable.

\subsection*{Formal Remark: Observational Coordinates}

Observed kinematic quantities correspond to integrals of entropy-weighted transport along admissible paths. Apparent acceleration reflects curvature in entropy descent trajectories when expressed in locally calibrated observational coordinates.

\section{Semantic and Cognitive Extensions}

Because the framework is fundamentally concerned with constraint, uncertainty, and flow, it extends naturally beyond physical systems. In semantic and cognitive domains, the base manifold $M$ is interpreted as a representational state space, $\Phi$ as uncertainty density, and $\vec{v}$ as directed inferential or attentional flow (Cover and Thomas, 2006; Landauer, 1961).

Cognitive stability emerges as the formation of semantic attractors: regions toward which uncertainty flow converges under admissible evolution. Concepts, beliefs, and skills are metastable configurations sustained by continuous entropy management rather than static storage. Learning corresponds to reshaping the constraint landscape itself through irreversible uncertainty reduction.

Agency is treated not as a primitive faculty but as a boundary condition on entropy flow. An agent is characterized by the capacity to modulate admissible vector fields, selectively directing entropy descent in accordance with goals or values. This does not collapse agency into regulation. Regulatory systems stabilize gradients locally; agents select which irreversible constraints to accept, thereby committing to particular futures.

Semantic equivalence arises through gauge structure. Distinct representations may be gauge-equivalent if they induce identical entropy flows and uncertainty reductions, even when differing in computational cost or accessibility. Meaning is invariant under representational transformation so long as constraint-respecting structure is preserved (Ellis, 2016).

\subsection*{Formal Perspective: Cognitive Entropy Balance}

Cognitive evolution obeys the same continuity relation
\[
\partial_t \Phi + \nabla \cdot (\Phi \vec{v}) = \sigma,
\]
with $\sigma \ge 0$ measuring irreversible information integration. Forgetting, revision, and learning are all constrained by entropy production bounds.

\section{AKSZ--BV Interpretation}

The coupled entropy–flow system admits a natural formulation within the AKSZ–BV framework, which encodes dynamics geometrically through graded symplectic structures. In this interpretation, the fields $(\Phi,\vec{v})$ arise as components of a mapping from a source manifold—representing spacetime or an abstract informational base—into a graded target space equipped with a shifted symplectic form (Baez and Biamonte, 2018).

The BV symplectic structure encodes admissible variations of entropy and flow, while gauge symmetries correspond to cohomological equivalences rather than physical degrees of freedom. The classical master equation enforces global consistency: entropy balance, conservation laws, and symmetry constraints are unified as a single nilpotency condition on the BV differential.

Time evolution is not fundamental in this picture. Apparent dynamics correspond to homological flow on the space of fields. Dissipation is represented as controlled deviation from exactness rather than as explicit forcing. The AKSZ–BV formalism thus provides a compact encoding of constraint-first dynamics, clarifying the inseparability of entropy, geometry, and symmetry.

\section{Meta-Theoretical Positioning and Methodological Commitments}

The entropy–flow framework is not offered as a competing fundamental theory, nor as a universal reductionist schema. It functions instead as an architectural language for organizing, comparing, and translating theories that are already successful within their domains. Its core commitment is that admissibility precedes dynamics: what matters first is what is allowed, not what happens.

This stance aligns partially with variational principles, statistical mechanics, and information geometry, while departing decisively from each. Like variational approaches, it encodes behavior through consistency conditions rather than explicit laws. Unlike classical mechanics, the action functional is not ontologically primary. Like statistical mechanics, it emphasizes entropy and irreversibility, but treats entropy as a field-level quantity governing structure at every scale (Jaynes, 2003; Prigogine, 1980).

The framework resists substance-based ontology. Fields, flows, and manifolds are not ontological furniture but elements of a bookkeeping system for constraint satisfaction. What is minimally real is the persistence of constraint-respecting structure. Objects, agents, and even spacetime itself are stabilized patterns within entropy flow rather than primitive constituents (Butterfield, 2014; Sklar, 1993).

This restraint enables cross-domain coherence without equivocation. The same formal relations apply across physics, cosmology, and cognition because the same architectural problem is being solved: how constrained systems redistribute degrees of freedom irreversibly while maintaining localized stability.

\section{Conclusion}

By integrating scalar entropy fields, vectorial transport, and geometric constraint formalisms, this essay has articulated a unified descriptive framework applicable across physical, cosmological, and cognitive domains. Its central claim is deliberately restrained: structure is the residual pattern left by entropy flow constrained by compatibility, conservation, and dissipation bounds.

Treating constraint as primary and geometry as expressive rather than prescriptive offers an alternative to substance-heavy and expansionist ontologies. The framework does not seek to replace existing theories but to reveal the shared architectural logic that underlies them, enabling translation, comparison, and synthesis across domains that are too often treated in isolation.

\appendix
\section{Entropy--Flow Field Equations}

Let $M$ be a smooth $n$-dimensional manifold with volume form $dV$.  
Let $\Phi : M \times \mathbb{R} \to \mathbb{R}_{\ge 0}$ be the entropy field and  
$\vec{v} : M \times \mathbb{R} \to TM$ the transport field.

The fundamental balance equation is
\begin{equation}
\partial_t \Phi + \nabla \cdot (\Phi \vec{v}) = \sigma,
\qquad
\sigma \ge 0.
\end{equation}

Assume a constitutive relation
\begin{equation}
\vec{v} = - \nabla \mu,
\end{equation}
where $\mu$ is an entropy potential (chemical potential analogue).

Define the entropy functional
\begin{equation}
\mathcal{E}[\Phi] = \int_M U(\Phi)\, dV,
\end{equation}
with $U''(\Phi) > 0$.

Then
\begin{equation}
\mu = \frac{\delta \mathcal{E}}{\delta \Phi} = U'(\Phi).
\end{equation}

Substituting yields the nonlinear diffusion equation
\begin{equation}
\partial_t \Phi
=
\nabla \cdot \left( \Phi \nabla U'(\Phi) \right) + \sigma.
\end{equation}

\subsection*{Entropy Production}

Define entropy production density
\begin{equation}
\sigma = \Phi \lVert \nabla \mu \rVert^2.
\end{equation}

Then
\begin{equation}
\frac{d}{dt} \mathcal{E}[\Phi]
=
- \int_M \Phi \lVert \nabla \mu \rVert^2 \, dV
\le 0.
\end{equation}

Thus $\mathcal{E}$ is a Lyapunov functional.

\subsection*{Wasserstein Gradient Flow}

Define the metric
\begin{equation}
\langle \xi, \eta \rangle_\Phi
=
\int_M \Phi \, \xi \cdot \eta \, dV.
\end{equation}

Then the evolution equation is the gradient flow
\begin{equation}
\partial_t \Phi
=
- \nabla_\Phi \mathcal{E},
\end{equation}
with respect to the Wasserstein-type metric induced by $\Phi$.

\subsection*{Worked Example: Porous Medium Case}

Let
\begin{equation}
U(\Phi) = \frac{1}{m-1} \Phi^m, \qquad m>1.
\end{equation}

Then
\begin{equation}
\mu = \Phi^{m-1},
\qquad
\vec{v} = - (m-1)\Phi^{m-2}\nabla \Phi.
\end{equation}

The evolution equation becomes
\begin{equation}
\partial_t \Phi
=
\nabla \cdot \left( \Phi^{m-1} \nabla \Phi \right),
\end{equation}
which is the porous medium equation.

Compactly supported initial data generate finite-speed propagation and self-similar attractors.
\section{Linear Stability and Spectral Structure}

Let $(\Phi_0,\vec{v}_0)$ be a stationary solution of the entropy--flow equations on $M$:
\begin{equation}
\partial_t \Phi_0 = 0,
\qquad
\nabla \cdot (\Phi_0 \vec{v}_0) = \sigma_0,
\qquad
\sigma_0 \ge 0.
\end{equation}

Assume $\vec{v}_0 = -\nabla \mu_0$ with $\mu_0 = U'(\Phi_0)$.

Introduce perturbations
\begin{equation}
\Phi = \Phi_0 + \varepsilon \,\delta\Phi,
\qquad
\mu = \mu_0 + \varepsilon \,\delta\mu,
\end{equation}
with
\begin{equation}
\delta\mu = U''(\Phi_0)\,\delta\Phi.
\end{equation}

\subsection*{Linearized Operator}

Linearizing
\[
\partial_t \Phi = \nabla\cdot(\Phi\nabla\mu)
\]
yields
\begin{equation}
\partial_t \delta\Phi
=
\mathcal{L}\,\delta\Phi,
\end{equation}
where
\begin{equation}
\mathcal{L}\delta\Phi
=
\nabla\cdot\!\left(
\Phi_0 \nabla (U''(\Phi_0)\delta\Phi)
+
\delta\Phi \nabla\mu_0
\right).
\end{equation}

If $\nabla\mu_0 = 0$ (equilibrium background), this simplifies to
\begin{equation}
\mathcal{L}\delta\Phi
=
\nabla\cdot\!\left(
\Phi_0 U''(\Phi_0)\nabla\delta\Phi
\right).
\end{equation}

\subsection*{Self-Adjoint Form}

Define weighted inner product
\begin{equation}
\langle f,g\rangle_{\Phi_0}
=
\int_M f g \,\Phi_0^{-1} \, dV.
\end{equation}

Then $\mathcal{L}$ is self-adjoint and non-positive:
\begin{equation}
\langle \delta\Phi, \mathcal{L}\delta\Phi\rangle_{\Phi_0}
=
-
\int_M
\Phi_0 U''(\Phi_0)
\lVert\nabla\delta\Phi\rVert^2
\, dV
\le 0.
\end{equation}

Hence all non-gauge modes decay.

\subsection*{Gauge Zero Modes}

Zero eigenvalues correspond to:
\begin{equation}
\nabla\delta\Phi = 0
\quad \Rightarrow \quad
\delta\Phi = \text{const}.
\end{equation}

These correspond to entropy-preserving relabelings and are factored out by gauge reduction.

\subsection*{Attractor Criterion}

A configuration $(\Phi_0,\vec{v}_0)$ is an attractor iff
\begin{equation}
\mathrm{spec}(\mathcal{L}) \subset (-\infty,0] \cup \{0_{\text{gauge}}\}.
\end{equation}

The basin of attraction is determined by nonlinear corrections
\begin{equation}
\partial_t \delta\Phi
=
\mathcal{L}\delta\Phi
+
\mathcal{N}(\delta\Phi),
\end{equation}
where $\mathcal{N}=O(\delta\Phi^2)$.


\subsection*{Worked Example: One-Dimensional Domain}

Let $M=[0,L]$ with periodic boundary conditions.
Let $\Phi_0 = \Phi_\star$ constant.

Then
\begin{equation}
\mathcal{L} = D\,\partial_x^2,
\qquad
D=\Phi_\star U''(\Phi_\star)>0.
\end{equation}

Eigenmodes:
\begin{equation}
\delta\Phi_k(x)=e^{ikx},
\qquad
\lambda_k = -Dk^2.
\end{equation}

All modes decay exponentially except $k=0$ (gauge).

Thus $\Phi_\star$ is globally stable.

\subsection*{Metastability}

If $U''(\Phi_0)$ vanishes at isolated points, $\mathcal{L}$ becomes degenerate.
This produces slow manifolds and long-lived metastable states.

Such degeneracies correspond to structural boundaries between attractors.

\section{Hamiltonian Reduction and Controlled Symplectic Breaking}

Let $(\Phi,\pi)$ denote the canonical field variables, where
\begin{equation}
\pi := \lambda
\end{equation}
is the momentum conjugate to the entropy field $\Phi$.

Define the Hamiltonian functional
\begin{equation}
\mathcal{H}[\Phi,\pi]
=
\int_M
\left(
\frac{1}{2}\,\Phi \lVert \nabla \pi \rVert^2
+
U(\Phi)
\right)\, dV.
\end{equation}

The canonical symplectic form on phase space is
\begin{equation}
\omega
=
\int_M d\pi \wedge d\Phi.
\end{equation}

Hamilton’s equations are
\begin{align}
\partial_t \Phi &= \frac{\delta \mathcal{H}}{\delta \pi}
= -\nabla\cdot(\Phi\nabla\pi), \\
\partial_t \pi &= -\frac{\delta \mathcal{H}}{\delta \Phi}
= -\frac{1}{2}\lVert\nabla\pi\rVert^2 - U'(\Phi).
\end{align}

\subsection*{Entropy Constraint as Phase-Space Restriction}

Entropy production is defined as
\begin{equation}
\sigma = \Phi \lVert \nabla \pi \rVert^2.
\end{equation}

The admissibility condition
\begin{equation}
\sigma \ge 0
\end{equation}
restricts phase-space trajectories to a cone
\begin{equation}
\mathcal{C} = \{ (\Phi,\pi) \mid \Phi \ge 0 \}.
\end{equation}

Phase-space volume contraction occurs whenever
\begin{equation}
\frac{d}{dt} \int_M \Phi \, dV > 0.
\end{equation}

Thus Liouville’s theorem fails globally.

\subsection*{Broken Symplectic Flow}

Let $X_\mathcal{H}$ be the Hamiltonian vector field.
Then
\begin{equation}
\mathcal{L}_{X_\mathcal{H}} \omega = 0
\end{equation}
only when $\sigma = 0$.

For $\sigma > 0$, the effective flow satisfies
\begin{equation}
\mathcal{L}_{X_\mathcal{H}} \omega = - d\Sigma,
\end{equation}
where $\Sigma$ is an entropy current one-form.

Thus symplecticity is broken in a controlled, monotone manner.

\subsection*{Reduced Dynamics}

Define the reduced phase space
\begin{equation}
\mathcal{P}_{\text{red}}
=
\{ (\Phi,\pi) \} / \sim
\end{equation}
where $\sim$ identifies trajectories differing by gauge-zero modes.

On $\mathcal{P}_{\text{red}}$, the effective dynamics is gradient-like:
\begin{equation}
\partial_t \Phi = - \nabla_\Phi \mathcal{E} + O(\sigma).
\end{equation}

Conservative Hamiltonian flow appears only on the boundary $\sigma = 0$.

\subsection*{Worked Example: Finite-Dimensional Reduction}

Let $\Phi \in \mathbb{R}_{>0}$ and $\pi \in \mathbb{R}$.

Hamiltonian:
\begin{equation}
H(\Phi,\pi) = \frac{1}{2}\Phi \pi^2 + U(\Phi).
\end{equation}

Equations:
\begin{align}
\dot{\Phi} &= -\Phi \pi, \\
\dot{\pi} &= -\frac{1}{2}\pi^2 - U'(\Phi).
\end{align}

Entropy production:
\begin{equation}
\sigma = \Phi \pi^2.
\end{equation}

Eliminate $\pi$ using $\pi = -\dot{\Phi}/\Phi$:
\begin{equation}
\dot{\Phi}^2 = 2\Phi\left(E - U(\Phi)\right),
\end{equation}
with $E$ non-increasing.

Thus the system is gradient-descent with inertial correction.

\subsection*{Irreversibility Criterion}

A trajectory is reversible iff
\begin{equation}
\sigma(t) \equiv 0
\quad\Rightarrow\quad
\pi(t)\equiv 0
\quad\Rightarrow\quad
\Phi = \text{const}.
\end{equation}

All nontrivial trajectories are irreversible.

\section{AKSZ--BV Formulation of Entropy--Flow Dynamics}

Let $\Sigma$ be a $d$-dimensional source manifold (spacetime or abstract base),
and let $\mathcal{M}$ be a graded target manifold equipped with a degree
$(d-1)$ symplectic form $\omega$.

Define coordinates on $\mathcal{M}$:
\begin{equation}
\Phi \in \mathcal{M}^{0},
\qquad
v_i \in \mathcal{M}^{-1},
\qquad
\pi \in \mathcal{M}^{d-2},
\end{equation}
where $\Phi$ is the entropy field, $v_i$ the transport field components,
and $\pi$ the BV momentum conjugate to $\Phi$.

\subsection*{BV Symplectic Structure}

Define the shifted symplectic form
\begin{equation}
\omega
=
\int_\Sigma
\delta \pi \wedge \delta \Phi
+
\delta v_i \wedge \delta v^{\dagger i},
\end{equation}
of total degree $-1$, where $v^{\dagger i}$ are antifields.

\subsection*{AKSZ Action}

Define the AKSZ action functional
\begin{equation}
S
=
\int_\Sigma
\left(
\pi \, d\Phi
+
\frac{1}{2}\Phi\, g^{ij} v_i v_j
-
U(\Phi)
\right),
\end{equation}
where $d$ is the de Rham differential on $\Sigma$.

Gauge symmetry:
\begin{equation}
\delta \Phi = \epsilon,
\qquad
\delta \pi = 0,
\qquad
\delta v_i = \partial_i \epsilon,
\end{equation}
corresponding to entropy-preserving relabelings.

\subsection*{Classical Master Equation}

The BV bracket satisfies
\begin{equation}
\{S,S\} = 0
\end{equation}
iff
\begin{equation}
\partial_t \Phi + \nabla \cdot (\Phi \vec{v}) = 0.
\end{equation}

Entropy production is introduced as controlled BV-breaking:
\begin{equation}
\{S,S\} = \Sigma,
\qquad
\Sigma \ge 0.
\end{equation}

Thus dissipation appears as homological non-closure.

\subsection*{Gauge Fixing}

Choose gauge-fixing fermion
\begin{equation}
\Psi = \int_\Sigma \pi\, \nabla \cdot v.
\end{equation}

Eliminating antifields yields
\begin{equation}
v_i = -\partial_i U'(\Phi),
\end{equation}
recovering gradient-flow dynamics.

\subsection*{Derived Critical Locus}

The derived critical locus of $S$ is the derived intersection
\begin{equation}
\mathrm{Crit}(S)
=
\{ \delta S = 0 \}
\end{equation}
encoding entropy balance, transport admissibility,
and gauge equivalence simultaneously.

\subsection*{Worked Toy Model: 1D Source}

Let $\Sigma = S^1$ with coordinate $x$.

Fields:
\[
\Phi(x), \quad v(x), \quad \pi(x).
\]

Action:
\begin{equation}
S
=
\int_{S^1}
\pi \partial_x \Phi
+
\frac{1}{2}\Phi v^2
-
U(\Phi)\, dx.
\end{equation}

Gauge-fixing gives
\begin{equation}
v = -\partial_x U'(\Phi),
\end{equation}
and reduced equation
\begin{equation}
\partial_t \Phi
=
\partial_x \big(\Phi\,\partial_x U'(\Phi)\big),
\end{equation}
recovering the entropy–flow PDE.

\subsection*{Interpretation}

Exact BV symmetry corresponds to reversible entropy transport.
Controlled BV anomaly encodes irreversible entropy production.

Homological flow replaces time evolution.

\section{Discrete Entropy--Flow Dynamics on a Lattice}

Let $M$ be discretized as a regular lattice with cells indexed by $i \in \mathbb{Z}^d$.
Let $\Phi_i^n \ge 0$ denote the cell-averaged entropy field at discrete time $t^n = n\Delta t$.

Define nearest-neighbor fluxes $F_{ij}^n$ across faces $(i,j)$.

\subsection*{Finite-Volume Update Rule}

The discrete entropy balance equation is
\begin{equation}
\Phi_i^{n+1}
=
\Phi_i^n
-
\frac{\Delta t}{\Delta x^d}
\sum_{j \sim i}
F_{ij}^n
+
\Delta t\,\sigma_i^n,
\end{equation}
where the sum runs over neighboring cells $j$.

Define fluxes by upwinded gradient descent:
\begin{equation}
F_{ij}^n
=
-\frac{1}{2}
\left(
\Phi_i^n + \Phi_j^n
\right)
\frac{U'(\Phi_j^n)-U'(\Phi_i^n)}{\Delta x}.
\end{equation}

\subsection*{Discrete Entropy Production}

Define local entropy production
\begin{equation}
\sigma_i^n
=
\sum_{j \sim i}
\frac{1}{2}
\left(
\Phi_i^n + \Phi_j^n
\right)
\left(
\frac{U'(\Phi_j^n)-U'(\Phi_i^n)}{\Delta x}
\right)^2
\ge 0.
\end{equation}

Thus positivity of entropy production is enforced by construction.

\subsection*{Discrete Lyapunov Functional}

Define discrete entropy functional
\begin{equation}
\mathcal{E}^n
=
\sum_i U(\Phi_i^n)\,\Delta x^d.
\end{equation}

Then
\begin{equation}
\mathcal{E}^{n+1} - \mathcal{E}^n
=
- \Delta t
\sum_i
\sigma_i^n
\Delta x^d
\le 0.
\end{equation}

Hence $\mathcal{E}^n$ is a discrete Lyapunov function.

\subsection*{CFL Condition}

Stability requires
\begin{equation}
\Delta t
\le
\frac{\Delta x^2}{\max_i \left(\Phi_i^n U''(\Phi_i^n)\right)}.
\end{equation}

This ensures monotone entropy descent and positivity preservation.

\subsection*{Worked Example: 1D Two-Cell System}

Let $i=1,2$, periodic boundary.

Initial data:
\begin{equation}
\Phi_1^0 = a, \qquad \Phi_2^0 = b.
\end{equation}

Let $U(\Phi)=\frac{1}{2}\Phi^2$.

Then
\begin{align}
F_{12}^0 &= -\frac{a+b}{2}(b-a), \\
\sigma_1^0 &= \sigma_2^0 = \frac{(a+b)}{2}(b-a)^2.
\end{align}

Update:
\begin{align}
\Phi_1^1 &= a + \Delta t (b-a), \\
\Phi_2^1 &= b - \Delta t (b-a).
\end{align}

Thus $\Phi_1^n,\Phi_2^n \to \frac{a+b}{2}$ monotonically.

\subsection*{Interpretation}

Discrete evolution converges to entropy-balanced attractors.
All numerical trajectories respect irreversibility by construction.
The scheme provides a faithful computational realization of
constraint-first entropy flow.

\section{Event-Historical Formulation of Entropy--Flow Dynamics}

Let $\mathcal{H}$ denote the space of admissible histories.
A history $h \in \mathcal{H}$ is a finite or countable sequence of events
\begin{equation}
h = (e_1, e_2, \dots, e_n),
\end{equation}
with each event $e_k$ corresponding to an irreversible update of the entropy field.

Let $\Phi^k$ denote the entropy configuration after event $e_k$.
Define the event update rule
\begin{equation}
\Phi^{k+1} = \mathcal{U}_{e_k}(\Phi^k),
\end{equation}
where $\mathcal{U}_{e_k}$ is a constraint-respecting transformation.

\subsection*{Admissibility and Partial Order}

Define entropy increment
\begin{equation}
\Delta S(e_k) := \mathcal{E}[\Phi^{k+1}] - \mathcal{E}[\Phi^k].
\end{equation}

Admissibility requires
\begin{equation}
\Delta S(e_k) \le 0.
\end{equation}

Define a partial order on histories:
\begin{equation}
h \preceq h'
\quad \text{iff} \quad
h' \text{ extends } h \text{ by admissible events}.
\end{equation}

This order replaces time as the primitive structure.
Temporal ordering is induced, not assumed.


\subsection*{Discrete Continuity Equation}

Let $i$ index lattice cells.
An event $e$ updates a finite neighborhood $N(e)$.

For $i \in N(e)$,
\begin{equation}
\Phi_i^{k+1}
=
\Phi_i^k
+
\sum_{j \sim i} J_{ij}(e)
+
\sigma_i(e),
\end{equation}
with
\begin{equation}
\sum_{i \in N(e)} \sigma_i(e) \ge 0,
\qquad
\sum_{i \in N(e)} J_{ij}(e) = 0.
\end{equation}

Thus each event locally conserves transport and globally produces entropy.

\subsection*{Event Composition}

Events compose via
\begin{equation}
\mathcal{U}_{e_{k+1}} \circ \mathcal{U}_{e_k},
\end{equation}
but composition is non-invertible whenever $\Delta S(e_k)>0$.

Define the event semigroup
\begin{equation}
\mathcal{E} = \langle e \mid \Delta S(e) \ge 0 \rangle,
\end{equation}
which replaces a group of time translations.

\subsection*{Continuum Limit}

Let events occur at rate $\lambda$.
Define
\begin{equation}
\Phi(x,t)
=
\lim_{\lambda \to \infty}
\Phi^{\lfloor \lambda t \rfloor}(x).
\end{equation}

Assuming bounded increments,
the discrete update converges to
\begin{equation}
\partial_t \Phi + \nabla \cdot (\Phi \vec{v}) = \sigma,
\end{equation}
recovering the entropy--flow PDE.

\subsection*{Worked Example: Two-Event System}

Let $\Phi \in \mathbb{R}^2$.

Event $e_1$:
\begin{equation}
\Phi \mapsto
\begin{pmatrix}
\Phi_1 - \alpha(\Phi_1-\Phi_2) \\
\Phi_2 + \alpha(\Phi_1-\Phi_2)
\end{pmatrix},
\qquad
\Delta S(e_1) = -\alpha(\Phi_1-\Phi_2)^2.
\end{equation}

Event $e_2$:
\begin{equation}
\Phi \mapsto
\begin{pmatrix}
\Phi_1 \\
\Phi_2
\end{pmatrix}
+
\beta
\begin{pmatrix}
1 \\
1
\end{pmatrix},
\qquad
\Delta S(e_2) = 2\beta.
\end{equation}

Admissibility requires $\beta \ge 0$.

Composition:
\begin{equation}
e_2 \circ e_1 \neq e_1 \circ e_2.
\end{equation}

Irreversibility arises from noncommutativity and entropy monotonicity.

\subsection*{Equivalence to Gradient Descent}

Define event potential
\begin{equation}
\Delta S(e) = -\langle \nabla \mathcal{E}, \delta\Phi_e \rangle.
\end{equation}

Admissible events satisfy
\begin{equation}
\delta\Phi_e = -\eta_e \nabla \mathcal{E},
\qquad
\eta_e > 0.
\end{equation}

Thus event histories are discrete gradient descent paths.

\subsection*{Interpretation}

Continuous time is a coarse-graining of irreversible event order.
Causality is defined by admissibility, not chronology.
The entropy functional orders histories.

\section{Categorical Semantics of Constraint-First Entropy Flow}

This appendix reformulates the entropy–flow framework in categorical terms.
No explanatory prose is provided beyond definitions and constructions.

\subsection*{Objects and Morphisms}

Let $\mathcal{C}$ be a category whose objects are admissible entropy configurations $\Phi$, understood as states of the entropy field that satisfy the global and local constraint conditions specified by the framework. A morphism in $\mathcal{C}$ is an irreversible, constraint-respecting transformation
\[
f : \Phi \to \Phi',
\]
which maps one admissible configuration to another in a manner compatible with entropy monotonicity. Admissibility of a morphism is characterized by the requirement
\[
\mathcal{E}(\Phi') \le \mathcal{E}(\Phi),
\]
where $\mathcal{E}$ denotes the entropy functional. This inequality expresses the irreversibility of morphisms in $\mathcal{C}$ and encodes the constraint that entropy cannot decrease along admissible transformations.

Composition in $\mathcal{C}$ is given by ordinary function composition and is associative by construction. Identity morphisms correspond to the trivial transformations for which $\Phi' = \Phi$ and the associated entropy production vanishes identically. As a result, $\mathcal{C}$ is generally non-groupoidal: morphisms need not be invertible, and invertibility occurs only in the degenerate case of zero entropy production.


\subsection*{Partial Order and Thin Category}

Define a preorder on objects by
\[
\Phi \preceq \Psi
\quad\Longleftrightarrow\quad
\exists\, f : \Phi \to \Psi.
\]

Entropy monotonicity implies antisymmetry up to gauge equivalence.
After quotienting by gauge, $\mathcal{C}$ is thin (at most one morphism between objects).

\subsection*{Monoidal Structure}

Define a symmetric monoidal product
\[
(\Phi_1,\Phi_2) \mapsto \Phi_1 \otimes \Phi_2
\]
with entropy functional
\[
\mathcal{E}(\Phi_1 \otimes \Phi_2)
=
\mathcal{E}(\Phi_1) + \mathcal{E}(\Phi_2).
\]

Monoidal unit is the minimal-entropy configuration $\Phi_\emptyset$.

Morphisms satisfy
\[
(f_1 \otimes f_2)(\Phi_1 \otimes \Phi_2)
=
f_1(\Phi_1)\otimes f_2(\Phi_2).
\]

\subsection*{Entropy as a Functor}

Define the entropy functor
\[
\mathcal{E} : \mathcal{C} \to (\mathbb{R}_{\le}, +),
\]
where $(\mathbb{R}_{\le},+)$ is the monoidal preorder category of real numbers
ordered by $\le$.

Functoriality:
\[
\mathcal{E}(g\circ f) \le \mathcal{E}(f),
\qquad
\mathcal{E}(f\otimes g)=\mathcal{E}(f)+\mathcal{E}(g).
\]

Thus entropy is a monoidal, order-preserving functor.

\subsection*{Gradient Flow as Enriched Morphism}

Let $\mathcal{C}$ be enriched over $(\mathbb{R}_{\ge 0},+)$.

Define enrichment by
\[
\mathrm{Hom}(\Phi,\Psi) = \inf \left\{
\int_0^T \sigma(t)\,dt
\;\middle|\;
\Phi \to \Psi
\right\}.
\]

Composition corresponds to addition of entropy cost:
\[
\mathrm{Hom}(\Phi,\Xi)
\le
\mathrm{Hom}(\Phi,\Psi)+\mathrm{Hom}(\Psi,\Xi).
\]

\subsection*{Gauge Equivalence as Groupoid Completion}

Let $\mathcal{G}$ denote the groupoid of gauge transformations.

Form the quotient category
\[
\mathcal{C}_{\text{phys}} = \mathcal{C}/\mathcal{G}.
\]

Objects are gauge-equivalence classes $[\Phi]$.
All morphisms respect gauge orbits.

$\mathcal{C}_{\text{phys}}$ is skeletal and thin.

\subsection*{Event Category}

We define an \emph{event category} $\mathcal{E}$ whose purpose is to encode irreversibility at the level of elementary transformations. The objects of $\mathcal{E}$ are entropy configurations $\Phi$, understood as admissible states of the entropy field. Morphisms in $\mathcal{E}$ are elementary events $e : \Phi \to \Phi'$, each representing a single irreversible update of the configuration. Admissibility of a morphism is determined by an entropy monotonicity condition: an event $e$ is permitted if and only if the associated entropy increment satisfies $\Delta S(e) \ge 0$. Composition of morphisms corresponds to temporal succession of events and is associative by construction, but generally non-invertible, reflecting the irreversible character of entropy-producing transformations.

There exists a canonical functor
\[
F : \mathcal{E} \to \mathcal{C}
\]
from the event category to the category $\mathcal{C}$ of admissible entropy transformations. This functor acts identically on objects and sends each elementary event to the induced constraint-respecting transformation on entropy configurations. Functoriality expresses the fact that sequential composition of events in $\mathcal{E}$ is mapped to composition of the corresponding irreversible morphisms in $\mathcal{C}$. In this way, $\mathcal{E}$ provides a generative, event-level description of entropy flow, while $\mathcal{C}$ captures its coarse-grained, configuration-level semantics.


Event composition is associative but non-invertible.
$\mathcal{E}$ is a semigroupoid.

\subsection*{Continuum Limit as Colimit}

Let $\{ \Phi^n \}_{n\in\mathbb{N}}$ be a directed system of events.

The continuum entropy flow is the colimit
\[
\Phi(t) = \varinjlim_{n\to\infty} \Phi^n.
\]

Dynamics arise as directed colimits of irreversible morphisms.

\subsection*{Adjunction: Constraint vs Dynamics}

Define categories:
\[
\mathcal{C}_{\text{adm}} \quad \text{(admissible configurations)},
\qquad
\mathcal{C}_{\text{dyn}} \quad \text{(dynamical evolutions)}.
\]

There exists an adjunction
\[
\mathcal{C}_{\mathrm{adm}}
\;\;\overset{L}{\underset{R}{\rightleftarrows}}\;\;
\mathcal{C}_{\mathrm{dyn}},
\]
between the category $\mathcal{C}_{\mathrm{adm}}$ of admissible configurations and the category $\mathcal{C}_{\mathrm{dyn}}$ of dynamical evolutions. The left adjoint $L$ assigns to each admissible configuration the class of entropy–gradient dynamics that are compatible with its constraint structure, thereby constructing a family of permissible flows from static admissibility data. Conversely, the right adjoint $R$ maps a given dynamical evolution to its underlying admissible configuration by forgetting the specific temporal realization while retaining the constraint relations that render the evolution permissible.

The existence of this adjunction formalizes the constraint-first methodology advocated throughout this work. Primacy is accorded to the right adjoint $R$, in the sense that admissibility is treated as conceptually prior to dynamics. Dynamical laws are not taken as fundamental objects but arise as left adjoint constructions from constraint data. In this way, the adjunction captures, in categorical form, the methodological inversion whereby the space of allowable configurations determines the space of possible evolutions, rather than the reverse.


\subsection*{Worked Example: Two-State Category}

Let $\Phi=\{\phi_1,\phi_2\}$ with
\[
\mathcal{E}(\phi_1) > \mathcal{E}(\phi_2).
\]

There exists a unique non-identity morphism
\[
f:\phi_1\to\phi_2,
\]
but no morphism $\phi_2\to\phi_1$.

Thus $\mathcal{C}$ is a poset category with terminal object $\phi_2$.

\subsection*{Interpretation}

Entropy defines a categorical order.
Irreversibility is functorial.
Dynamics are directed morphisms.
Time is the nerve of a thin monoidal category.


\begin{thebibliography}{99}

\bibitem{Jaynes1957}
E.~T.~Jaynes,
Information theory and statistical mechanics,
\emph{Physical Review}, 106(4):620--630, 1957.

\bibitem{Jaynes2003}
E.~T.~Jaynes,
\emph{Probability Theory: The Logic of Science},
Cambridge University Press, 2003.

\bibitem{Prigogine1980}
I.~Prigogine,
\emph{From Being to Becoming},
W.~H.~Freeman, 1980.

\bibitem{PrigogineStengers1984}
I.~Prigogine and I.~Stengers,
\emph{Order Out of Chaos},
Bantam Books, 1984.

\bibitem{Onsager1931}
L.~Onsager,
Reciprocal relations in irreversible processes,
\emph{Physical Review}, 37:405--426, 1931.

\bibitem{Callen1985}
H.~B.~Callen,
\emph{Thermodynamics and an Introduction to Thermostatistics},
Wiley, 1985.

\bibitem{CoverThomas2006}
T.~M.~Cover and J.~A.~Thomas,
\emph{Elements of Information Theory},
Wiley, 2006.

\bibitem{Villani2009}
C.~Villani,
\emph{Optimal Transport: Old and New},
Springer, 2009.

\bibitem{Otto2001}
F.~Otto,
The geometry of dissipative evolution equations,
\emph{CPDE}, 26:101--174, 2001.

\bibitem{Arnold1989}
V.~I.~Arnold,
\emph{Mathematical Methods of Classical Mechanics},
Springer, 1989.

\bibitem{BaezBiamonte2018}
J.~C.~Baez and J.~Biamonte,
Quantum techniques for stochastic mechanics,
World Scientific, 2018.

\bibitem{BaezPollard2017}
J.~C.~Baez and B.~Pollard,
A compositional framework for reaction networks,
\emph{Rev. Math. Phys.}, 2017.

\bibitem{Gromov1999}
M.~Gromov,
\emph{Metric Structures},
Birkhäuser, 1999.

\bibitem{Ellis2016}
G.~F.~R.~Ellis,
How can physics underlie the mind?,
\emph{Topoi}, 35:247--260, 2016.

\bibitem{Butterfield2014}
J.~Butterfield,
Reduction, emergence, and renormalization,
\emph{Journal of Philosophy}, 111:5--49, 2014.

\bibitem{Sklar1993}
L.~Sklar,
\emph{Physics and Chance},
Cambridge University Press, 1993.

\bibitem{Landauer1961}
R.~Landauer,
Irreversibility and heat generation in the computing process,
\emph{IBM J. Res. Dev.}, 5:183--191, 1961.

\end{thebibliography}

\end{document}
